\Needspace{8\baselineskip}
\Article{292}
We distinguish decompositions of numbers into three squares (as we did
above for binary forms) from representations by the form $xx+yy+zz$ by
considering only the sizes of the squares in the former, but also their
order and the signs of their roots in the latter. So we regard the
representations $x=a$, $y=b$, $z=c$ and $x=a'$, $y=b'$, $z=c'$ as
distinct unless $a=a'$, $b=b'$, $c=c'$ all hold. But we regard the
decompositions into $aa+bb+cc$ and into $a'a'+b'b'+c'c'$ as one if these squares
are equal to those without regard to order. It follows that

I.\ A decomposition of the number $M$ into the squares $aa+bb+cc$
is equivalent to $48$ representations if none is $=0$ and
all are unequal; but to $24$ if \textit{either} one is $=0$
and the others unequal, \textit{or} none is $=0$ and two are equal.
If, however, in a decomposition of a given number into three squares,
two of them are $=0$, \textit{or} one is $=0$ and the others
equal, \textit{or} all are equal, it will be equivalent to $6$,
\textit{or} $12$, \textit{or} $8$ representations, respectively.
But these can occur only in the exceptional cases $M=1$, $2$,
$3$, respectively, since the representations must be proper.
Excluding these, let the number of all decompositions of the number
$M$ into three squares (without a common divisor) be $E$,
and suppose that among them there are $e$ with one square $0$
and $e'$ with two equal squares. The former can also be regarded
as decompositions into two squares, and the latter as decompositions
into a square and twice a square. Then the number of all proper
representations of the number $M$ by $xx+yy+zz$ will be
\[
=24(e+e')+48(E-e-e')=48E-24(e+e')
\]
\SourceLeaf{351}
But from the theory of binary forms we readily deduce that $e$
will be either $=0$ or $=2^{\mu-1}$, according as $-1$ is a quadratic
nonresidue or residue of $M$, and also $e'=0$ or $=2^{\mu-1}$, according
as $-2$ is a nonresidue or residue of $M$. Here $\mu$ denotes
the number of (odd) prime factors of $M$ (see Article 182; we
omit a fuller account here). It follows readily that
\[
\begin{aligned}
E&=2^{\mu-2}k,&&\text{if both $-1$ and $-2$ are $N.R.$ of $M$;}\\
E&=2^{\mu-2}(k+2),&&\text{if both numbers are residues; finally}\\
E&=2^{\mu-2}(k+1),&&\text{if one is a residue and the other a nonresidue.}
\end{aligned}
\]
In the excluded cases $M=1$ and $M=2$, this formula would give
$E=\tfrac{3}{4}$, whereas it should be $E=1$. But for $M=3$ it correctly
gives $E=1$, the exceptions compensating for each other.

So, if $M$ is a prime, $\mu=1$, and therefore $E=\tfrac{1}{2}(k+2)$ when
$M\equiv1\ (\mathrm{mod.}\ 8)$; $E=\tfrac{1}{2}(k+1)$ when $M\equiv3$ or $\equiv5$. These special theorems
were discovered by induction by the distinguished Le Gendre and given
in his excellent paper, already cited several times, \textit{Hist.
de l'Ac. de Paris} 1785, p.\ 530 ff., though in a somewhat different
form. This is mainly because he did not distinguish proper from
improper equivalence, and so combined opposite classes.

II.\ To find all decompositions of the number $M$ into three
squares (without a common divisor), we do not need to find all proper
representations of all the forms $\varphi$, $\varphi'$, $\varphi''$. First,
it is easy to confirm that all $(48)$ representations of the form
$\varphi$ belonging to the same value of the expression $\surd{}-(p,-q,r)$ (setting
$\varphi=(p,q,r)$) give the same decomposition of the number $M$. So it
is enough to have one of them, or, what comes to the same thing, only
to write down all distinct decompositions\HistoricalNote{*)}{We must
always understand \textit{proper}, if we choose to transfer this
term from representations to decompositions.} of the form $\varphi$
into three squares; and likewise for the others $\varphi'$, $\varphi''$,
and so on. Next, if $\varphi$ belongs to a nonambiguous class, we may
entirely omit the form chosen from the opposite class; in other words,
it is enough to consider just one of two opposite classes. For the
choice of a form from each class is completely arbitrary. Suppose
that from the class opposite to the one containing $\varphi$ we choose
the form opposite to $\varphi$, and let it be $=\varphi'$. Then we readily
see that, if
\SourceLeaf{352}
the proper decompositions of the form \(\varphi\) are expressed generally by
\[
(gt+hu)^2+(g't+h'u)^2+(g''t+h''u)^2
\]
all decompositions of the form \(\varphi'\) will be expressed by
\[
(gt-hu)^2+(g't-h'u)^2+(g''t-h''u)^2
\]
and that these give the same decompositions of the number \(M\)
as the former. Finally, if \(\varphi\) is a form in an ambiguous class,
but is neither in the principal class nor equivalent to the form
\((2,0,\frac12 M)\) or \((2,1,\frac12(M+1))\) (according as \(M\) is even or odd),
we may omit half the values of the expression \(\surd-(p,-q,r)\). For brevity,
however, we do not discuss this shortcut further here. --- We can
also use these same shortcuts when all proper representations of
\(M\) by \(xx+yy+zz\) are wanted, since these are very easily obtained
from the decompositions.

As an example, we will find all decompositions of the number \(770\)
into three squares. Here \(\mu=3\), \(e=e'=0\), and therefore \(E=2k\).
Classifying the positive binary forms of determinant \(-770\) gives
the number of positive classes as \(=32\). We omit this classification
for brevity, since anyone can easily construct it by the method of
Article 231. All these classes are properly primitive and distributed
among \(8\) genera, so that \(k=4\), and therefore \(E=8\). The
genus whose characteristic number is \(-1\) must clearly have the
particular characters with respect to the numbers \(5,7,11\) given by
\(R5;\ N7;\ N11\).
From Article 263, we readily conclude that its character with respect
to the number \(8\) must be \(1\text{ and }3,\ 8\). Now in the genus with character
\(1\text{ and }3,\ 8;\ R5;\ N7;\ N11\) there are four classes, as whose representatives we choose
the forms \((6,2,129)\), \((6,-2,129)\), \((19,3,41)\), \(19,-3,41)\). But we discard the
second and fourth classes, as they are opposite to the first and third.
We have already given four decompositions of the form \((19,3,41)\) in
Article 289. These give decompositions of the number \(770\) into
\(9+361+400\), \(16+25+729\), \(81+400+289\), \(576+169+25\). In the same way we find four
decompositions of the form \(6tt+4tu+129uu\) into
\[
\begin{gathered}
(t-8u)^2+(2t+u)^2+(t+8u)^2,\\
(t-10u)^2+(2t+5u)^2+(t+2u)^2\\
(2t-5u)^2+(t+10u)^2+(t+2u)^2,\\
(2t+7u)^2+(t-8u)^2+(t-4u)^2
\end{gathered}
\]
\SourceLeaf{353}
arising, respectively, from the following values of the expression
\(\surd-(6,-2,129)\): \((48,369)\), \((62,-149)\), \((92,-159)\), \((202,61)\). These give decompositions
of the number \(770\) into \(225+256+289\), \(1+144+625\), \(64+81+625\), \(16+225+529\).
Besides these eight decompositions there are no others.

What concerns decompositions of numbers into three squares having
common divisors follows so easily from the general theory of Article
281 that there is no need to spend time on it.

\Needspace{8\baselineskip}
\begin{center}\textit{Proof of Fermat's theorems that every integer can be decomposed into three triangular numbers or four squares.}\end{center}
\Article{293}
The preceding investigations also give a proof of the famous theorem
that \textit{every positive integer can be decomposed into three
triangular numbers}, discovered long ago by Fermat, but for which a
rigorous proof had been missing until now. Clearly any decomposition
of the number \(M\) into triangular numbers
\[
\tfrac12 x(x+1)+\tfrac12 y(y+1)+\tfrac12 z(z+1)
\]
gives a decomposition of the number \(8M+3\) into three odd squares
\[
(2x+1)^2+(2y+1)^2+(2z+1)^2
\]
and conversely. But by the preceding theory every positive integer
\(8M+3\) can be resolved into three squares, which must be odd (see
the note to Article 291). The number of such resolutions depends
both on the number of prime factors of \(8M+3\) and on the number
of classes into which the binary forms of determinant \(-(8M+3)\) are
distributed. There will be just as many decompositions of the number
\(M\) into three triangular numbers. We assume, however, that
\(\frac12 x(x+1)\) is regarded as a triangular number for any integer value
of \(x\). If one preferred to exclude zero, the theorem would
have to be changed to this: Every positive integer is either itself
a triangular number or can be resolved into two or three triangular
numbers. A similar change would be needed in the following theorem
if one preferred to exclude zero from the squares.

The same principles prove another theorem of Fermat: \textit{every
positive integer can be decomposed into four squares.} Subtract
any square (smaller than the number) from a number of the form
\(4n+2\), an even square from a number of the form
\SourceLeaf{354}
\(4n+1\), or an odd square from a number of the form \(4n+3\). In
all these cases the remainder can be resolved into three squares,
and so the given number into four. Finally, a number of the form
\(4n\) can be expressed as \(4^{\mu}N\), where \(N\) belongs to
one of the preceding three forms. Resolving \(N\) into four
squares also resolves \(4^{\mu}N\). From a number of the form \(8n+3\)
we can also subtract a square whose root is divisible by four; from
a number of the form \(8n+7\), a square whose root is twice an odd
number; and from a number of the form \(8n+4\), an odd square.
The remainder can then be resolved into three squares. This theorem
had already been proved by the distinguished La Grange, \textit{Nouv.
Mém. de l'Ac. de Berlin} 1770, p.\ 123. His proof, entirely different
from ours, was explained more fully by the distinguished Euler in
\textit{Actis Ac. Petr. Vol.}\ II, p.\ 48. --- Other theorems of Fermat,
which form a kind of continuation of these, stating that every integer
can be resolved into five pentagonal numbers, six hexagonal numbers,
seven heptagonal numbers, \textit{and so on}, still lack a proof
and seem to require deeper principles.
